\section[Aksesibilitas: Prinsip dan Praktik]{Aksesibilitas (a11y): Prinsip dan Praktik}

Aksesibilitas web (sering disingkat \textbf{a11y}---``a'' diikuti 11 huruf dan ``y'') adalah praktik memastikan bahwa situs web dapat digunakan oleh \textbf{semua orang}, termasuk pengguna dengan disabilitas visual, auditori, motorik, atau kognitif. Menurut WHO, sekitar 16\% populasi dunia hidup dengan disabilitas, sehingga aksesibilitas bukan sekadar fitur tambahan melainkan kebutuhan fundamental. Di banyak negara, aksesibilitas web bahkan merupakan kewajiban hukum. Bagi pengembang web, membangun situs yang aksesibel juga berarti menjangkau audiens yang lebih luas dan meningkatkan kualitas keseluruhan produk \cite{mdnAccessibility}.

Pedoman utama untuk aksesibilitas web adalah \textbf{Web Content Accessibility Guidelines (WCAG)}, yang dikembangkan oleh W3C. WCAG dibangun di atas empat prinsip dasar yang dikenal dengan akronim \textbf{POUR}: Perceivable (dapat dipersepsi), Operable (dapat dioperasikan), Understandable (dapat dipahami), dan Robust (tangguh). Setiap prinsip memiliki kriteria sukses yang dikelompokkan dalam tiga level konformitas: A (minimum), AA (standar yang direkomendasikan), dan AAA (optimal) \cite{wcagW3c}.

\begin{figure}[htbp]
\centering
\resizebox{0.85\textwidth}{!}{%
\begin{tikzpicture}[node distance=1.8cm, transform shape, every node/.style={font=\footnotesize}]
  \node (pour) [class, minimum width=4cm, fill=blue!25] {\textbf{POUR}\\Prinsip Aksesibilitas};

  \node (p) [class, below left=1.5cm and 2.5cm of pour, fill=orange!15, minimum width=3.5cm, minimum height=1.5cm]
    {\textbf{Perceivable}\\Konten dapat\\dipersepsi indra};
  \node (o) [class, below left=1.5cm and -0.5cm of pour, fill=green!15, minimum width=3.5cm, minimum height=1.5cm]
    {\textbf{Operable}\\Antarmuka dapat\\dioperasikan};
  \node (u) [class, below right=1.5cm and -0.5cm of pour, fill=yellow!15, minimum width=3.5cm, minimum height=1.5cm]
    {\textbf{Understandable}\\Konten dapat\\dipahami};
  \node (r) [class, below right=1.5cm and 2.5cm of pour, fill=purple!15, minimum width=3.5cm, minimum height=1.5cm]
    {\textbf{Robust}\\Kompatibel dengan\\berbagai teknologi};

  \draw [arrow] (pour) -- (p);
  \draw [arrow] (pour) -- (o);
  \draw [arrow] (pour) -- (u);
  \draw [arrow] (pour) -- (r);
\end{tikzpicture}%
}
\caption{Empat prinsip aksesibilitas POUR menurut WCAG}
\end{figure}

\begin{table}[htbp]
\centering
\begin{tabular}{|P{1.5cm}|P{3.5cm}|P{3.5cm}|P{3.5cm}|}
\hline
\textbf{Level} & \textbf{Deskripsi} & \textbf{Contoh Kriteria} & \textbf{Rekomendasi} \\
\hline
A & Kebutuhan dasar & Teks alt pada gambar; konten navigable via keyboard & Wajib dipenuhi \\
\hline
AA & Standar industri & Kontras warna minimum 4.5:1; resize teks hingga 200\% & Target utama \\
\hline
AAA & Level optimal & Kontras 7:1; bahasa isyarat untuk media & Ideal namun tidak selalu feasible \\
\hline
\end{tabular}
\caption{Level konformitas WCAG dan contoh kriteria sukses}
\end{table}

Dalam praktik HTML, aksesibilitas dimulai dengan penggunaan elemen semantik yang tepat. Elemen-elemen seperti \code{<nav>}, \code{<main>}, \code{<header>}, \code{<footer>}, \code{<button>}, dan \code{<a>} secara otomatis menyediakan informasi peran (role) dan perilaku (behavior) kepada teknologi bantu. Atribut \code{alt} pada \code{<img>} memberikan deskripsi teks untuk pengguna yang tidak dapat melihat gambar. Navigasi keyboard harus berfungsi untuk semua elemen interaktif, dan urutan fokus (tab order) harus logis \cite{webdevAccessibility}.

Ketika elemen semantik HTML tidak cukup untuk menyampaikan makna antarmuka, atribut \textbf{ARIA} (Accessible Rich Internet Applications) dapat digunakan. ARIA menyediakan tiga jenis atribut: \code{role} (mendefinisikan peran elemen, misalnya \code{role="dialog"}), \code{aria-*} states (status, misalnya \code{aria-expanded="true"}), dan \code{aria-*} properties (properti, misalnya \code{aria-label="Menu utama"}). Namun, aturan pertama ARIA adalah: \textbf{``Jangan gunakan ARIA jika Anda bisa menggunakan elemen HTML semantik yang sudah memiliki makna dan perilaku bawaan''} \cite{mdnAccessibility}.

\begin{webcode}[language=HTML, caption={Contoh penggunaan ARIA untuk navigasi dan tombol}]
<!-- Nav dengan aria-label untuk membedakan navigasi -->
<nav aria-label="Menu Utama">
  <ul>
    <li><a href="/">Beranda</a></li>
    <li><a href="/produk">Produk</a></li>
    <li><a href="/kontak">Kontak</a></li>
  </ul>
</nav>

<!-- Tombol toggle dengan ARIA state -->
<button aria-expanded="false"
        aria-controls="panel-info">
  Tampilkan Info
</button>
<div id="panel-info" hidden>
  <p>Informasi tambahan di sini.</p>
</div>

<!-- Gambar dekoratif (diabaikan screen reader) -->
<img src="dekorasi.png" alt="" aria-hidden="true" />

<!-- Gambar informatif (deskripsi wajib) -->
<img src="grafik-penjualan.png"
     alt="Grafik penjualan naik 25% di Q3 2026" />
\end{webcode}

\begin{konsep}
\textbf{Aturan Pertama ARIA:} Sebelum menggunakan atribut ARIA, tanyakan: ``Apakah ada elemen HTML semantik yang sudah menyediakan makna ini?'' Misalnya:
\begin{itemize}
  \item Gunakan \code{<button>} alih-alih \code{<div role="button">}.
  \item Gunakan \code{<nav>} alih-alih \code{<div role="navigation">}.
  \item Gunakan \code{<input type="checkbox">} alih-alih \code{<span role="checkbox">}.
\end{itemize}
ARIA tidak mengubah perilaku elemen; ia hanya mengubah \textbf{informasi yang disampaikan ke teknologi bantu}. Elemen \code{<div role="button">} tidak akan bisa diklik atau difokuskan dengan keyboard kecuali Anda menambahkan \code{tabindex} dan event handler secara manual. Elemen \code{<button>} sudah menyediakan semua ini secara bawaan \cite{webdevAccessibility}.
\end{konsep}

